% GENERATED FILE. Edit canonical lessons, then run npm run generate:books.
% canonical-source-hash: fnv1a64:bfdc13b8e6fbe6bf
% canonical-lessons: TE-C194-intipani, TE-C194-panimanisi, TE-C194-darji, TE-C194-vyapari, TE-C194-nanem

\chapter{Housework, Domestic helper, Tailor, Trader, Coin}
\label{ch:te-housework-domestic-helper-tailor-trader-coin}
\hlchaptermodality{194}

\begin{chapteropening}
\textbf{By the end of this chapter:} \emph{I can say housework, a domestic helper, a tailor, a trader and a coin.}

Complete the last lesson of chapter 194: I can say housework, a domestic helper, a tailor, a trader and a coin.
\end{chapteropening}

\section[\texorpdfstring{\te{ఇంటిపని} (iṇṭipani)}{ఇంటిపని (iṇṭipani)}]{\te{ఇంటిపని} (iṇṭipani) — housework}

\label{lesson:TE-C194-intipani}



\begin{quote}
Before the new one: say the Telugu for a triangle, then the Telugu for a wheel.
\end{quote}

\subsection*{You'll want to know: \te{ఇంటిపని}}
\textbf{\te{ఇంటిపని}} — \emph{iṇṭipani} — \textquotedblleft{}housework\textquotedblright{}.

\begin{grammarlens}[title={What you've built}]
One more word for this chapter.
\end{grammarlens}

\subsection*{Your turn}
\begin{itemize}
\raggedright
  \item \emph{Say it:} \emph{iṇṭipani}
  \item \emph{Say it:} \emph{iṇṭipani}, once more
  \item \emph{Say it:} say \emph{iṇṭipani}
  \item \emph{Recall:} say the Telugu for a triangle, then the Telugu for a wheel, then say \emph{iṇṭipani} again
\end{itemize}

\begin{tcolorbox}[breakable,colback=teal!4,colframe=teal!35!black,title={Before you move on}]
What does \te{ఇంటిపని} mean? (\textquotedblleft{}Housework\textquotedblright{}.) Say it once more.
\end{tcolorbox}
\section[\texorpdfstring{\te{పనిమనిషి} (panimaniṣi)}{పనిమనిషి (panimaniṣi)}]{\te{పనిమనిషి} (panimaniṣi) — a domestic helper}

\label{lesson:TE-C194-panimanisi}



\begin{quote}
Before the new one: say the Telugu for a wheel, then the Telugu for housework.
\end{quote}

\subsection*{You'll want to know: \te{పనిమనిషి}}
\textbf{\te{పనిమనిషి}} — \emph{panimaniṣi} — \textquotedblleft{}a domestic helper\textquotedblright{}.

\begin{grammarlens}[title={What you've built}]
One more word for this chapter.
\end{grammarlens}

\subsection*{Your turn}
\begin{itemize}
\raggedright
  \item \emph{Say it:} \emph{panimaniṣi}
  \item \emph{Say it:} \emph{panimaniṣi}, once more
  \item \emph{Say it:} say \emph{panimaniṣi}
  \item \emph{Recall:} say the Telugu for a wheel, then the Telugu for housework, then say \emph{panimaniṣi} again
\end{itemize}

\begin{tcolorbox}[breakable,colback=teal!4,colframe=teal!35!black,title={Before you move on}]
What does \te{పనిమనిషి} mean? (\textquotedblleft{}A domestic helper\textquotedblright{}.) Say it once more.
\end{tcolorbox}
\section[\texorpdfstring{\te{దర్జీ} (darjī)}{దర్జీ (darjī)}]{\te{దర్జీ} (darjī) — a tailor}

\label{lesson:TE-C194-darji}



\begin{quote}
Before the new one: say the Telugu for housework, then the Telugu for a domestic helper.
\end{quote}

\subsection*{You'll want to know: \te{దర్జీ}}
\textbf{\te{దర్జీ}} — \emph{darjī} — \textquotedblleft{}a tailor\textquotedblright{}.

\begin{grammarlens}[title={What you've built}]
One more word for this chapter.
\end{grammarlens}

\subsection*{Your turn}
\begin{itemize}
\raggedright
  \item \emph{Say it:} \emph{darjī}
  \item \emph{Say it:} \emph{darjī}, once more
  \item \emph{Say it:} say \emph{darjī}
  \item \emph{Recall:} say the Telugu for housework, then the Telugu for a domestic helper, then say \emph{darjī} again
\end{itemize}

\begin{tcolorbox}[breakable,colback=teal!4,colframe=teal!35!black,title={Before you move on}]
What does \te{దర్జీ} mean? (\textquotedblleft{}A tailor\textquotedblright{}.) Say it once more.
\end{tcolorbox}
\section[\texorpdfstring{\te{వ్యాపారి} (vyāpāri)}{వ్యాపారి (vyāpāri)}]{\te{వ్యాపారి} (vyāpāri) — a trader, a merchant}

\label{lesson:TE-C194-vyapari}



\begin{quote}
Before the new one: say the Telugu for a domestic helper, then the Telugu for a tailor.
\end{quote}

\subsection*{You'll want to know: \te{వ్యాపారి}}
\textbf{\te{వ్యాపారి}} — \emph{vyāpāri} — \textquotedblleft{}a trader, a merchant\textquotedblright{}.

\begin{grammarlens}[title={What you've built}]
One more word for this chapter.
\end{grammarlens}

\subsection*{Your turn}
\begin{itemize}
\raggedright
  \item \emph{Say it:} \emph{vyāpāri}
  \item \emph{Say it:} \emph{vyāpāri}, once more
  \item \emph{Say it:} say \emph{vyāpāri}
  \item \emph{Recall:} say the Telugu for a domestic helper, then the Telugu for a tailor, then say \emph{vyāpāri} again
\end{itemize}

\begin{tcolorbox}[breakable,colback=teal!4,colframe=teal!35!black,title={Before you move on}]
What does \te{వ్యాపారి} mean? (\textquotedblleft{}A trader, a merchant\textquotedblright{}.) Say it once more.
\end{tcolorbox}
\section[\texorpdfstring{\te{నాణెం} (nāṇeṁ)}{నాణెం (nāṇeṁ)}]{\te{నాణెం} (nāṇeṁ) — a coin}

\label{lesson:TE-C194-nanem}



\begin{quote}
Before the new one: say the Telugu for a tailor, then the Telugu for a trader.
\end{quote}

\subsection*{You'll want to know: \te{నాణెం}}
\textbf{\te{నాణెం}} — \emph{nāṇeṁ} — \textquotedblleft{}a coin\textquotedblright{}.

\begin{grammarlens}[title={What you've built}]
One more word for this chapter.
\end{grammarlens}

\subsection*{Your turn}
\begin{itemize}
\raggedright
  \item \emph{Say it:} \emph{nāṇeṁ}
  \item \emph{Say it:} \emph{nāṇeṁ}, once more
  \item \emph{Say it:} say \emph{nāṇeṁ}
  \item \emph{Recall:} say the Telugu for a tailor, then the Telugu for a trader, then say \emph{nāṇeṁ} again
\end{itemize}

\begin{tcolorbox}[breakable,colback=teal!4,colframe=teal!35!black,title={Before you move on}]
What does \te{నాణెం} mean? (\textquotedblleft{}A coin\textquotedblright{}.) Say it once more.
\end{tcolorbox}
